\chapter{Foundations of Downstream Separation}
\label{ch:sep}

\begin{quote}\itshape
``Don't waste clean thinking on dirty enzymes.'' --- Efraim Racker, on why
purity comes first.
\end{quote}

\noindent
The models of Chapter~\ref{ch:models} speak fluently of isotherms, plates,
gradients, and breakthrough; the fermentation of Chapter~\ref{ch:ferment} of
substrate, biomass, and specific growth rates. To a statistician or a software
engineer drawn to this monograph by its design methods, that vocabulary can be a
wall. This chapter dismantles it. We build, from the molecule up, the small set
of physical ideas a reader needs to follow the rest of the book: what a
\emph{downstream process} is and why it exists, what makes biological molecules
hard to separate, how an \emph{adsorption isotherm} governs everything, how a
\emph{column} turns that isotherm into a separation, and what a
\emph{fermentation} is as a unit operation. As in Chapter~\ref{ch:stats}, the
treatment is self-contained but curated to exactly what Chapters~\ref{ch:models}
through~\ref{ch:ferment} use, and the quantitative figures are generated from a
small teaching module (\code{vlab\_doe.foundations.separation\_demo}).

%==============================================================================
\section{The bioprocess and its downstream half}
\label{sec:sep-train}

A biological drug---a monoclonal antibody, an enzyme, a vaccine antigen---is made
in two great phases. \emph{Upstream}, living cells (bacterial, yeast, or
mammalian) are grown in a bioreactor and coaxed into making the target molecule;
the output is a turbid, dilute broth in which the product is a small fraction of
a vast chemical zoo. \emph{Downstream} is everything after: the sequence of
\emph{unit operations}---discrete, repeatable processing steps---that clarifies,
concentrates, and purifies that broth into a pure drug substance
(Figure~\ref{fig:sep-train}). The chromatography models at the heart of
Part~II live in this downstream half, though the monograph ranges beyond
it---upstream into fermentation, and out to the benchtop.

A platform downstream train for an antibody runs: \emph{harvest and
clarification} (centrifuge and filter away the cells and debris); \emph{capture}
(a highly selective first chromatography step---classically Protein~A---that binds
the product and lets most impurities wash through); one or two \emph{polishing}
steps (further chromatography, usually ion exchange and/or hydrophobic
interaction, that remove the last, most product-like impurities); and
\emph{formulation} by ultrafiltration/diafiltration (UF/DF), which concentrates
the product and exchanges it into its final buffer. Purity climbs at every
stage.

Three numbers grade a step, and they are the responses optimised throughout the
book. \emph{Yield} (or recovery) is the fraction of product that survives the
step. \emph{Purity} is the product's share of what comes out. \emph{Productivity}
is throughput---product mass per unit resin per unit time. They almost always
\emph{compete}: collecting a wider, purer cut sacrifices yield; running faster
sacrifices resolution. Navigating that trade-off is what the experimental
designs of Chapter~\ref{ch:design} are for.

\begin{figure}[t]
\centering
\includegraphics[width=\linewidth]{foundations_bioprocess.png}
\caption[The bioprocess train]{A platform bioprocess. \emph{Upstream}
(fermentation/cell culture) makes the molecule; \emph{downstream}---the focus of
Chapters~\ref{ch:models} and~\ref{ch:design}---clarifies, captures, polishes, and
formulates it. Each box is a
\emph{unit operation}, and purity increases from left to right. The
chromatographic steps (capture and polish) are modelled in
Chapter~\ref{ch:models}; the fermentation in Chapter~\ref{ch:ferment}.}
\label{fig:sep-train}
\end{figure}

%==============================================================================
\section{The molecules, and why they are hard to separate}
\label{sec:sep-molecules}

A protein is a long chain of \emph{amino acids} folded into a compact shape. Two
of its surface properties do almost all the work of separation. The first is
\emph{charge}: several amino-acid side chains gain or lose protons as the
solution pH changes, so a protein's \emph{net charge} swings from positive at
low pH to negative at high pH, passing through zero at a characteristic pH called
the \emph{isoelectric point} (pI). The second is \emph{hydrophobicity}: patches
of oily, water-avoiding residues on the surface. A given protein is thus a unique
point in a space of size, charge-versus-pH, and hydrophobicity---and each axis is
the basis of a different separation \emph{mode} (Section~\ref{sec:sep-modes}).

The difficulty is that the impurities resemble the product. The broth carries
\emph{host-cell proteins} (hundreds of the cell's own proteins), \emph{DNA},
\emph{endotoxin}, leached ligand from a previous column, and---most insidiously
---\emph{product-related} variants: aggregates (clumps of the product), fragments,
and charge variants that differ from the molecule we want by the smallest of
margins. Separating a protein from buffer is trivial; separating it from a
slightly larger clump of itself is the entire art, and it is why the field needs
several orthogonal modes and the careful optimisation this book describes.

\begin{appbox}{why orthogonality matters}
Because impurities span charge, size, and hydrophobicity, no single step removes
them all. The platform pairs \emph{orthogonal} modes---ion exchange (charge) then
hydrophobic interaction (hydrophobicity)---so that an impurity which co-elutes
with the product on one basis is resolved on the other. Choosing and sequencing
these modes is a downstream-development decision the virtual laboratory lets one
rehearse \emph{in silico}.
\end{appbox}

%==============================================================================
\section{Equilibrium and the adsorption isotherm}
\label{sec:sep-isotherm}

Almost every purification step rests on \emph{partitioning}: when a solute is
offered a choice between the flowing liquid (the \emph{mobile phase}) and a solid
(the \emph{stationary phase}, or \emph{adsorbent}/\emph{resin}---typically porous
beads packed in a tube), it distributes between the two according to how strongly
it is attracted to the solid. At equilibrium the relationship between the amount
on the solid, $q$, and the amount in the liquid, $c$, is the \emph{adsorption
isotherm}---the single most important constitutive law in the book.

At low concentration the isotherm is \emph{linear},
%
\begin{equation}
q = H\,c,
\label{eq:sep-henry}
\end{equation}
%
where the slope $H$, the \emph{Henry constant}, is the partition coefficient: a
large $H$ means the solute strongly prefers the solid and so moves slowly through
a column. But the solid has only finitely many binding sites, so as $c$ rises the
loading must saturate. The simplest law capturing both regimes is the
\emph{Langmuir isotherm} (Figure~\ref{fig:sep-isotherm}),
%
\begin{equation}
q = \frac{q_{\max}\,b\,c}{1 + b\,c},
\label{eq:sep-langmuir}
\end{equation}
%
with $q_{\max}$ the \emph{saturation capacity} (the monolayer ceiling) and $b$ an
affinity constant. Its two limits are exactly the two regimes of preparative
chromatography. When $bc\ll1$ the denominator is one and $q\approx(q_{\max}b)\,c$:
the linear law~\eqref{eq:sep-henry} with $H=q_{\max}b$, valid for a dilute,
analytical injection. When $bc\gg1$, $q\to q_{\max}$: the column is
\emph{overloaded}, as a preparative step deliberately is to maximise
productivity. The bend between them is what lets a strongly bound species
displace a weaker one---the nonlinear effect that makes high-resolution
preparative separation possible, and the reason Chapter~\ref{ch:models} builds
its whole engine on a multi-component Langmuir core.

\begin{keybox}{Algorithm: the isotherm, and where $H$ hides in it}
\code{separation\_demo.langmuir(c, q\_max, b)} is a one-liner,
\begin{verbatim}
    q = q_max * b * c / (1.0 + b * c)
\end{verbatim}
and the Henry constant is recovered as its initial slope, $H=q_{\max}b$
(\code{separation\_demo.henry}). Every chromatography mode in
Chapter~\ref{ch:models} is this \emph{same} expression; the modes differ only in
how the affinity $b$ responds to a mobile-phase ``modulator'' such as salt
(Section~\ref{sec:sep-modes}).
\end{keybox}

\begin{figure}[t]
\centering
\includegraphics[width=0.72\linewidth]{foundations_isotherm.png}
\caption[The adsorption isotherm]{The Langmuir isotherm (blue) and its dilute,
linear limit (red dashed). At low concentration the loading follows the Henry
line of slope $H=q_{\max}b$; as concentration rises it bends over toward the
finite saturation capacity $q_{\max}$. Where a band sits is set by $H$; how a
loaded column behaves is set by the bend. Generated by
\code{separation\_demo.langmuir}.}
\label{fig:sep-isotherm}
\end{figure}

%==============================================================================
\section{The column: retention, plates, and resolution}
\label{sec:sep-column}

A chromatography \emph{column} is a tube packed with adsorbent beads through
which mobile phase flows. The natural unit of throughput is the \emph{column
volume} (CV), the liquid volume of the packed bed; the fraction of that volume
occupied by flowing liquid (between and within the beads) is the \emph{porosity}.
Because a retained solute spends part of its time stuck to the solid, it travels
more slowly than the liquid: it emerges after a \emph{retention volume} larger
than one CV, by a factor set directly by the isotherm slope $H$. The
\emph{retention factor} $k$ measures this extra holdup, and a more strongly bound
solute (larger $H$, larger $k$) elutes later. This is the lever every separation
pulls: arrange conditions so that two solutes have \emph{different} retention,
and they come off the column at different times.

\paragraph{Breakthrough and capacity.}
If feed is pumped continuously onto a column, the adsorbent fills from the inlet
and eventually the solute ``breaks through'' at the outlet---the
\emph{breakthrough curve} (Figure~\ref{fig:sep-column}, left). The simplest
faithful model imagines the column as a chain of $N$ well-mixed equilibrium
stages, the \emph{theoretical plates}; the outlet concentration then rises along
a sigmoid whose sharpness grows with $N$. This is the operative picture of a
\emph{capture} step run in \emph{bind-and-elute} mode (load until just before
breakthrough, then elute the captured product in a small, concentrated volume),
and---inverted---of \emph{flow-through} polishing, where the product passes
straight through while impurities bind.

\paragraph{Plates, peaks, and resolution.}
The same plate count governs how \emph{sharp} an eluting band is. A column of $N$
plates produces an approximately Gaussian peak whose width relative to its
retention time shrinks as $\sqrt N$: formally $N=t_R^2/\sigma^2$, so more plates
mean a taller, narrower peak---higher \emph{efficiency}. The practical question
is whether two neighbouring peaks are actually \emph{separated}, quantified by
the \emph{resolution}
%
\begin{equation}
R_s = \frac{2\,|t_{R,2}-t_{R,1}|}{w_1+w_2},
\qquad w = 4\sigma,
\label{eq:sep-rs}
\end{equation}
%
the distance between peak centres in units of their average width
(Figure~\ref{fig:sep-column}, right). The universal target is $R_s\gtrsim1.5$,
``baseline'' separation, at which two equal peaks barely overlap. Resolution thus
has two ingredients: \emph{selectivity} (getting the retention times apart,
through the isotherm and the gradient) and \emph{efficiency} (keeping the peaks
narrow, through the plate count). Chapter~\ref{ch:models} computes exactly these
quantities from its simulated chromatograms.

\begin{figure}[t]
\centering
\includegraphics[width=\linewidth]{foundations_column.png}
\caption[Plates, breakthrough, and resolution]{\emph{Left:} breakthrough curves
from the tanks-in-series model---more theoretical plates $N$ give a sharper front,
so plate count is column efficiency made quantitative. \emph{Right:} two Gaussian
bands of the same efficiency; their separation is summarised by the resolution
$R_s$ of Eq.~\eqref{eq:sep-rs}, with $R_s\gtrsim1.5$ the usual target. Generated
by \code{separation\_demo.tanks\_in\_series\_breakthrough}, \code{gaussian\_peak},
and \code{resolution}.}
\label{fig:sep-column}
\end{figure}

%==============================================================================
\section{Why peaks broaden: convection, dispersion, and mass transfer}
\label{sec:sep-masstransfer}

In an ideal column every molecule of a solute would travel at the same speed and
emerge as an infinitely sharp spike. Real bands spread, and understanding why is
the bridge to the transport models of Chapter~\ref{ch:models}. Three mechanisms
broaden a band. \emph{Convection} simply carries the solute downstream at the
mobile-phase velocity---it moves the band but does not broaden it. \emph{Axial
dispersion} does broaden it: molecules take tortuous, unequal paths through the
packing and spread out, much as a diffusing ink drop blurs. And \emph{mass
transfer resistance}---the finite time a molecule needs to diffuse out of the
flowing stream, across the stagnant film around a bead, and into the bead's
pores to reach a binding site---means the solid is never quite in equilibrium with
the passing liquid, which smears the band further.

A model can either \emph{lump} all of this into a single effective dispersion and
a single mass-transfer rate constant---fast, and adequate for band positions and
widths---or \emph{resolve} the film and pore-diffusion steps explicitly, which is
slower but mechanistically faithful. That is exactly the choice
Chapter~\ref{ch:models} makes between its lumped ``transport-dispersive'' engine
and its ``general rate model.'' The vocabulary of this section---convection,
dispersion, film, pore diffusion---is all one needs to read that distinction.

%==============================================================================
\section{The modes in one idea: a modulator controls binding}
\label{sec:sep-modes}

The different chromatography \emph{modes} exploit the different surface
properties of Section~\ref{sec:sep-molecules}, but they share a single
organising idea: a knob in the mobile phase---the \emph{modulator}---sets how
strongly the solute binds, and a \emph{gradient} in that knob elutes the bound
solutes in order. The three modes the book treats differ only in what the
modulator is and which way binding responds:

\begin{itemize}
  \item \emph{Ion-exchange} (IEX) separates on \emph{charge}. The modulator is
  \emph{salt}: at low salt the charged protein binds the oppositely charged
  resin; raising the salt screens the attraction and elutes it. A second knob,
  \emph{pH}, tunes the protein's net charge relative to its pI.
  \item \emph{Hydrophobic-interaction} (HIC) separates on \emph{hydrophobicity}.
  The modulator is again salt, but with the \emph{opposite} sense: high salt
  drives the protein onto the hydrophobic resin (``salting out''), and
  \emph{lowering} the salt elutes it.
  \item \emph{Reversed-phase} (RP) also separates on hydrophobicity, with the
  modulator the fraction of \emph{organic solvent}; raising it elutes.
\end{itemize}
%
Two everyday measures of the salt knob recur: \emph{ionic strength}, the total
concentration of dissolved ions, and \emph{conductivity}, its easily measured
electrical proxy. A \emph{buffer} holds the pH steady while the salt or organic
gradient does the eluting. The remarkable economy that Chapter~\ref{ch:models}
exploits is that all three modes are the \emph{same} Langmuir
isotherm~\eqref{eq:sep-langmuir}; only the law by which the affinity $b$ depends
on the modulator changes from one mode to the next. Understanding that single
sentence is most of what the modelling chapter assumes.

%==============================================================================
\section{The living unit operation: fermentation basics}
\label{sec:sep-fermentation}

The last unit operation in this book is upstream and alive: a
\emph{fermentation}, in which microbes convert a feedstock into a product---the
subject of Chapter~\ref{ch:ferment}, where the product is cultured milk. A few
ideas make that chapter's models legible. The microbes consume a
\emph{substrate} (a nutrient such as the milk sugar lactose) and grow as
\emph{biomass}; the rate at which biomass increases per unit biomass already
present is the \emph{specific growth rate} $\mu$. In a sealed \emph{batch}, $\mu$
is not constant: growth passes through a \emph{lag} phase (cells adjusting to the
new environment), an \emph{exponential} phase (unrestricted growth), a
\emph{deceleration} as the substrate runs low, and a \emph{stationary} plateau
when it is exhausted (Figure~\ref{fig:sep-growth}).

The workhorse description of the substrate's effect is \emph{Monod's law},
%
\begin{equation}
\mu = \mu_{\max}\,\frac{S}{K_S+S},
\label{eq:sep-monod}
\end{equation}
%
a saturating dependence: growth is near its maximum $\mu_{\max}$ while substrate
$S$ is plentiful and throttles toward zero as $S$ falls below the
half-saturation constant $K_S$. Biomass and substrate are tied together by a
\emph{yield} coefficient $Y_{X/S}$, the biomass formed per unit substrate
consumed. The few extra ingredients Chapter~\ref{ch:ferment} adds---a mechanistic
lag, a temperature optimum, inhibition by the acid the microbes excrete, and
cooperation between strains---are all refinements of this skeleton. Two final
terms: the small dose of culture used to start a batch is the \emph{inoculum},
and a defined mixture of microbial \emph{strains} used to inoculate is a
\emph{starter culture}---and choosing that mixture from a large library is the
combinatorial design problem that closes the book.

\begin{figure}[t]
\centering
\includegraphics[width=0.72\linewidth]{foundations_growth.png}
\caption[Batch microbial growth]{A simulated batch fermentation. Biomass (green)
moves through lag, exponential, deceleration, and stationary phases as the
substrate (amber, dashed) is consumed; the coupling is Monod's saturating
kinetics, Eq.~\eqref{eq:sep-monod}, with a yield linking the two. This skeleton
is what the richer model of Chapter~\ref{ch:ferment} elaborates. Generated by
\code{separation\_demo.monod\_batch}.}
\label{fig:sep-growth}
\end{figure}

\begin{keybox}{The thread into the next chapters}
With this chapter in hand, the rest of the book reads as a sequence of natural
elaborations. The \emph{isotherm} (Section~\ref{sec:sep-isotherm}) is the
constitutive heart of every model in Chapter~\ref{ch:models}; \emph{retention,
plates, and resolution} (Section~\ref{sec:sep-column}) are the quantities those
models compute; \emph{dispersion and mass transfer}
(Section~\ref{sec:sep-masstransfer}) distinguish its two solvers; the
\emph{modulator} idea (Section~\ref{sec:sep-modes}) unifies its three modes; and
\emph{Monod growth} (Section~\ref{sec:sep-fermentation}) is the seed of the
living reactor of Chapter~\ref{ch:ferment}. The physics supplies the questions;
the statistics of Chapter~\ref{ch:stats} and the designs of
Chapter~\ref{ch:design} supply the means to answer them efficiently.
\end{keybox}
